% ========================================================================================
% CAPITOLO 4 - COMUNICAZIONE CON HANDSHAKING
% ========================================================================================
\chapter{Comunicazione con handshaking}

% ----------------------------------------------------------------------------------------
\section{Esercizio 8: Comunicazione con handshaking}
% ----------------------------------------------------------------------------------------

\subsection{Esercizio 8.1}

\begin{traccia}
	Progettare, implementare in VHDL e testare mediante simulazione un sistema composto da 2
	nodi, A e B, che comunicano mediante un protocollo di handshaking. Il nodo A e il nodo B
	possiedono entrambi una memoria interna in cui sono memorizzate N stringhe di M bit,
	denominate X(i) e Y(i) rispettivamente (i=0,..,N-1). Il nodo A trasmette a B ciascuna
	stringa X(i) utilizzando un protocollo di handshaking; B, ricevuta la stringa X(i), calcola
	S(i)=X(i)+Y(i) e immagazzina la somma in opportune locazioni della propria memoria interna.

	Per il progetto è possibile considerare una implementazione di tipo comportamentale per
	effettuare la somma, mentre è necessario prevedere esplicitamente un componente contatore
	sia nel sistema A sia nel sistema B per scandire la trasmissione/ricezione delle stringhe e
	per terminare la comunicazione.
\end{traccia}

\progetto

\begin{scelta}[Handshaking a quattro fasi]
	I nodi comunicano con un protocollo a quattro fasi sui segnali \sig{req} (da A a B) e
	\sig{ack} (da B ad A), mostrato in Figura~\ref{fig:hs4}.
	\todo{Motivare la scelta rispetto al protocollo a due fasi.}
\end{scelta}

\begin{figure}[H]
	\centering
	\begin{tikzpicture}[yscale=0.55, xscale=1.25]
		% griglia delle fasi
		\foreach \x/\n in {1/1, 2.5/2, 5/3, 6.5/4} {
			\draw[grigliatempo, dashed] (\x,0.6) -- (\x,-5.4);
			\node[circle, draw, inner sep=1pt, font=\scriptsize] at (\x,1.3) {\n};
		}
		% dato
		\node[left] at (-0.1,-0.5) {\sig{dato}};
		\draw[onda] (0,0) -- (0.8,0) -- (1,-1) -- (5.3,-1) -- (5.5,0) -- (8.5,0);
		\draw[onda] (0,-1) -- (0.8,-1) -- (1,0) -- (5.3,0) -- (5.5,-1) -- (8.5,-1);
		\node[font=\small] at (3.15,-0.5) {$X(i)$ valido};
		% req
		\node[left] at (-0.1,-2.5) {\sig{req}};
		\draw[onda] (0,-3) -- (1,-3) -- (1,-2) -- (5,-2) -- (5,-3) -- (8.5,-3);
		% ack
		\node[left] at (-0.1,-4.5) {\sig{ack}};
		\draw[onda] (0,-5) -- (2.5,-5) -- (2.5,-4) -- (6.5,-4) -- (6.5,-5) -- (8.5,-5);
		% relazioni causa-effetto
		\draw[->, attcol, thick] (1.05,-2.1) to[bend left=15] (2.45,-4.1);
		\draw[->, attcol, thick] (2.55,-3.9) to[bend left=15] (4.95,-2.2);
		\draw[->, attcol, thick] (5.05,-2.9) to[bend right=15] (6.45,-4.1);
	\end{tikzpicture}
	\caption{Handshaking a quattro fasi. \fase{1} A presenta $X(i)$ e alza \sig{req};
	\fase{2} B cattura il dato, salva $S(i)$ e alza \sig{ack};
	\fase{3} A abbassa \sig{req}; \fase{4} B abbassa \sig{ack}.}
	\label{fig:hs4}
\end{figure}

\todo{Schema a blocchi dei due nodi (parte operativa e di controllo), automi delle unità di controllo.}

\implementazione
\todo{Codice dei nodi A e B e del sistema.}

\simulazione
\todo{Verifica di tutte le $S(i)$; prova con un nodo B più lento.}
